% !TeX program = lualatex
% !TeX root = pred9.tex
% ============================================================
% Zajednički preamble za sva predavanja RWA kolegija
% Uključiti s: % ============================================================
% Zajednički preamble za sva predavanja RWA kolegija
% Uključiti s: % ============================================================
% Zajednički preamble za sva predavanja RWA kolegija
% Uključiti s: \input{../common/preamble.tex}
% ============================================================

\documentclass[aspectratio=149,11pt]{beamer}
\usetheme{CambridgeUS}

\usepackage{fontspec}              % radi samo sa XeLaTeX/LuaLaTeX
\setsansfont{Latin Modern Sans}    % sans-serif font
\setmonofont{Latin Modern Mono}    % monospace font za kod
\usepackage[croatian]{babel}
\usepackage{csquotes}              % \enquote{} → hrvatski navodnici „..."
\usepackage{lmodern}
\usepackage{booktabs}
\usepackage{listings}
\usepackage{xcolor}
\usepackage{hyperref}

% ---------- Autor i logo ----------
\author{mag.ing. Josip Torić (SIT UNIZD)}
\titlegraphic{\includegraphics[height=2.5cm]{../common/SIT_logo_tekst1.pdf}}

% ---------- Listings ----------
\definecolor{codebg}{RGB}{248,248,248}
\definecolor{codeframeclr}{RGB}{220,220,220}
\lstdefinelanguage{bash}{
  morekeywords={sudo,apt,brew,pip,pipx,uv,winget,python,git,flask,source,deactivate,code,gunicorn,docker,compose,heroku,render,railway,node,npm,pnpm,yarn,uvicorn,alembic,pytest,ruff,black,precommit},
  sensitive=false,
  morecomment=[l]{\#}
}
\lstdefinelanguage{python}{
  morekeywords={from,import,def,class,return,if,elif,else,for,in,while,try,except,raise,with,as,not,and,or,is,lambda,yield,True,False,None,await,async},
  sensitive=true,
  morecomment=[l]{\#},
  morestring=[b]',
  morestring=[b]"
}
\lstdefinelanguage{sql}{
  morekeywords={SELECT,FROM,WHERE,INSERT,INTO,VALUES,UPDATE,SET,DELETE,CREATE,TABLE,PRIMARY,KEY,FOREIGN,REFERENCES,UNIQUE,NOT,NULL,INDEX,ON,ALTER,DROP,CONSTRAINT},
  sensitive=false,
  morecomment=[l]{--}
}
\lstdefinelanguage{json}{
  morecomment=[l]{//},
  morestring=[b]",
  showstringspaces=false
}
\lstdefinelanguage{typescript}{
  morekeywords={import,export,from,const,let,var,function,return,if,else,for,of,in,while,switch,case,break,async,await,try,catch,finally,throw,new,class,this,extends,implements,interface,type,enum,readonly,public,private,protected,static,get,set,true,false,null,undefined,void,string,number,boolean,any,unknown,never,typeof,keyof,as},
  sensitive=true,
  morecomment=[l]{//},
  morecomment=[s]{/*}{*/},
  morestring=[b]',
  morestring=[b]",
  morestring=[b]`
}
\lstalias{javascript}{typescript}
\lstset{
  backgroundcolor=\color{codebg},
  basicstyle=\ttfamily\footnotesize,
  frame=single,
  rulecolor=\color{codeframeclr},
  frameround=tttt,
  breaklines=true,
  tabsize=2,
  showstringspaces=false,
  upquote=true,
  columns=fullflexible
}

% ---------- Helper macros ----------
\newcommand{\codefile}[1]{\texttt{#1}}
\newcommand{\kbd}[1]{\texttt{#1}}

% ---------- Code frame environment ----------
% Automatski dodaje fragile i shrink=5, koristi se umjesto \begin{frame}
% Primjer: \begin{codeframe}{Naslov} ... \end{codeframe}
% Primjer s opcijama: \begin{codeframe}[shrink=15]{Naslov} ... \end{codeframe}
\newenvironment{codeframe}[2][]
  {\begin{frame}[fragile,shrink=5,#1]{#2}}
  {\end{frame}}

% ============================================================

\documentclass[aspectratio=149,11pt]{beamer}
\usetheme{CambridgeUS}

\usepackage{fontspec}              % radi samo sa XeLaTeX/LuaLaTeX
\setsansfont{Latin Modern Sans}    % sans-serif font
\setmonofont{Latin Modern Mono}    % monospace font za kod
\usepackage[croatian]{babel}
\usepackage{csquotes}              % \enquote{} → hrvatski navodnici „..."
\usepackage{lmodern}
\usepackage{booktabs}
\usepackage{listings}
\usepackage{xcolor}
\usepackage{hyperref}

% ---------- Autor i logo ----------
\author{mag.ing. Josip Torić (SIT UNIZD)}
\titlegraphic{\includegraphics[height=2.5cm]{../common/SIT_logo_tekst1.pdf}}

% ---------- Listings ----------
\definecolor{codebg}{RGB}{248,248,248}
\definecolor{codeframeclr}{RGB}{220,220,220}
\lstdefinelanguage{bash}{
  morekeywords={sudo,apt,brew,pip,pipx,uv,winget,python,git,flask,source,deactivate,code,gunicorn,docker,compose,heroku,render,railway,node,npm,pnpm,yarn,uvicorn,alembic,pytest,ruff,black,precommit},
  sensitive=false,
  morecomment=[l]{\#}
}
\lstdefinelanguage{python}{
  morekeywords={from,import,def,class,return,if,elif,else,for,in,while,try,except,raise,with,as,not,and,or,is,lambda,yield,True,False,None,await,async},
  sensitive=true,
  morecomment=[l]{\#},
  morestring=[b]',
  morestring=[b]"
}
\lstdefinelanguage{sql}{
  morekeywords={SELECT,FROM,WHERE,INSERT,INTO,VALUES,UPDATE,SET,DELETE,CREATE,TABLE,PRIMARY,KEY,FOREIGN,REFERENCES,UNIQUE,NOT,NULL,INDEX,ON,ALTER,DROP,CONSTRAINT},
  sensitive=false,
  morecomment=[l]{--}
}
\lstdefinelanguage{json}{
  morecomment=[l]{//},
  morestring=[b]",
  showstringspaces=false
}
\lstdefinelanguage{typescript}{
  morekeywords={import,export,from,const,let,var,function,return,if,else,for,of,in,while,switch,case,break,async,await,try,catch,finally,throw,new,class,this,extends,implements,interface,type,enum,readonly,public,private,protected,static,get,set,true,false,null,undefined,void,string,number,boolean,any,unknown,never,typeof,keyof,as},
  sensitive=true,
  morecomment=[l]{//},
  morecomment=[s]{/*}{*/},
  morestring=[b]',
  morestring=[b]",
  morestring=[b]`
}
\lstalias{javascript}{typescript}
\lstset{
  backgroundcolor=\color{codebg},
  basicstyle=\ttfamily\footnotesize,
  frame=single,
  rulecolor=\color{codeframeclr},
  frameround=tttt,
  breaklines=true,
  tabsize=2,
  showstringspaces=false,
  upquote=true,
  columns=fullflexible
}

% ---------- Helper macros ----------
\newcommand{\codefile}[1]{\texttt{#1}}
\newcommand{\kbd}[1]{\texttt{#1}}

% ---------- Code frame environment ----------
% Automatski dodaje fragile i shrink=5, koristi se umjesto \begin{frame}
% Primjer: \begin{codeframe}{Naslov} ... \end{codeframe}
% Primjer s opcijama: \begin{codeframe}[shrink=15]{Naslov} ... \end{codeframe}
\newenvironment{codeframe}[2][]
  {\begin{frame}[fragile,shrink=5,#1]{#2}}
  {\end{frame}}

% ============================================================

\documentclass[aspectratio=149,11pt]{beamer}
\usetheme{CambridgeUS}

\usepackage{fontspec}              % radi samo sa XeLaTeX/LuaLaTeX
\setsansfont{Latin Modern Sans}    % sans-serif font
\setmonofont{Latin Modern Mono}    % monospace font za kod
\usepackage[croatian]{babel}
\usepackage{csquotes}              % \enquote{} → hrvatski navodnici „..."
\usepackage{lmodern}
\usepackage{booktabs}
\usepackage{listings}
\usepackage{xcolor}
\usepackage{hyperref}

% ---------- Autor i logo ----------
\author{mag.ing. Josip Torić (SIT UNIZD)}
\titlegraphic{\includegraphics[height=2.5cm]{../common/SIT_logo_tekst1.pdf}}

% ---------- Listings ----------
\definecolor{codebg}{RGB}{248,248,248}
\definecolor{codeframeclr}{RGB}{220,220,220}
\lstdefinelanguage{bash}{
  morekeywords={sudo,apt,brew,pip,pipx,uv,winget,python,git,flask,source,deactivate,code,gunicorn,docker,compose,heroku,render,railway,node,npm,pnpm,yarn,uvicorn,alembic,pytest,ruff,black,precommit},
  sensitive=false,
  morecomment=[l]{\#}
}
\lstdefinelanguage{python}{
  morekeywords={from,import,def,class,return,if,elif,else,for,in,while,try,except,raise,with,as,not,and,or,is,lambda,yield,True,False,None,await,async},
  sensitive=true,
  morecomment=[l]{\#},
  morestring=[b]',
  morestring=[b]"
}
\lstdefinelanguage{sql}{
  morekeywords={SELECT,FROM,WHERE,INSERT,INTO,VALUES,UPDATE,SET,DELETE,CREATE,TABLE,PRIMARY,KEY,FOREIGN,REFERENCES,UNIQUE,NOT,NULL,INDEX,ON,ALTER,DROP,CONSTRAINT},
  sensitive=false,
  morecomment=[l]{--}
}
\lstdefinelanguage{json}{
  morecomment=[l]{//},
  morestring=[b]",
  showstringspaces=false
}
\lstdefinelanguage{typescript}{
  morekeywords={import,export,from,const,let,var,function,return,if,else,for,of,in,while,switch,case,break,async,await,try,catch,finally,throw,new,class,this,extends,implements,interface,type,enum,readonly,public,private,protected,static,get,set,true,false,null,undefined,void,string,number,boolean,any,unknown,never,typeof,keyof,as},
  sensitive=true,
  morecomment=[l]{//},
  morecomment=[s]{/*}{*/},
  morestring=[b]',
  morestring=[b]",
  morestring=[b]`
}
\lstalias{javascript}{typescript}
\lstset{
  backgroundcolor=\color{codebg},
  basicstyle=\ttfamily\footnotesize,
  frame=single,
  rulecolor=\color{codeframeclr},
  frameround=tttt,
  breaklines=true,
  tabsize=2,
  showstringspaces=false,
  upquote=true,
  columns=fullflexible
}

% ---------- Helper macros ----------
\newcommand{\codefile}[1]{\texttt{#1}}
\newcommand{\kbd}[1]{\texttt{#1}}

% ---------- Code frame environment ----------
% Automatski dodaje fragile i shrink=5, koristi se umjesto \begin{frame}
% Primjer: \begin{codeframe}{Naslov} ... \end{codeframe}
% Primjer s opcijama: \begin{codeframe}[shrink=15]{Naslov} ... \end{codeframe}
\newenvironment{codeframe}[2][]
  {\begin{frame}[fragile,shrink=5,#1]{#2}}
  {\end{frame}}


\title[Razvoj web aplikacija]{Razvoj web aplikacija - Predavanje 9 \\ State management, UX stanja i forme}
\date{07.12.2026.}

\begin{document}

\frame{\titlepage}

% ------------------------------------------------------------
\begin{frame}{Plan predavanja}
\begin{enumerate}
  \item State management: zašto i kad koristiti store
  \item Pinia store za poslovne entitete
  \item Lista s loading/error/empty stanjima
  \item Composables: ponovno iskoristiva reaktivna logika
  \item Forme s klijentskom i serverskom validacijom
  \item Registracije i status badgevi
  \item Reusable komponente (tablica, forma)
  \item Commitovi za ovo predavanje
\end{enumerate}

\vspace{0.6em}
Na kraju predavanja: club korisnik može vidjeti svoje natjecatelje, dodati novog i prijaviti ga na natjecanje.
\end{frame}

% ------------------------------------------------------------
\begin{frame}{Ishodi učenja}
Nakon ovog predavanja student može:
\begin{itemize}
  \item kreirati Pinia store za CRUD operacije nad entitetom
  \item implementirati loading, error i empty stanja u viewu
  \item izdvojiti reaktivnu logiku u composable funkciju
  \item napraviti formu s \texttt{v-model} bindingom i klijentskom validacijom
  \item prikazati serverske greške u formi za statuse 422 i 409
  \item izdvojiti ponavljajući UI u reusable komponentu s tipiziranim propsima
\end{itemize}
\end{frame}

% ------------------------------------------------------------
\begin{frame}{Kad koristiti store, a kad lokalni state}
\begin{tabular}{@{}lp{5cm}p{5cm}@{}}
  \toprule
  & \textbf{Lokalni state (\texttt{ref()})} & \textbf{Pinia store} \\
  \midrule
  Scope & Jedna komponenta & Cijela aplikacija \\
  Primjer & Form inputi, toggle & Lista natjecatelja, auth \\
  Persistencija & Do unmounta & Do refresha (ili localStorage) \\
  Dijeljenje & Props/emit & Import u bilo kojoj komponenti \\
  \bottomrule
\end{tabular}

\vspace{0.6em}
Pravilo: ako dva ili više viewova trebaju iste podatke, stavi u store. Ako je podatak relevantan samo unutar jedne komponente, drži ga lokalno.
\end{frame}

% ------------------------------------------------------------
\begin{frame}{State management: problem koji rješava}
Bez centralnog storea, dijeljenje podataka među komponentama zahtijeva \textbf{prop drilling}:

\vspace{0.4em}
\begin{center}
\texttt{App $\rightarrow$ Layout $\rightarrow$ Sidebar $\rightarrow$ UserInfo}
\end{center}

\vspace{0.4em}
Svaka komponenta u lancu mora primiti i proslijediti prop koji koristi samo zadnja. Ovo je:
\begin{itemize}
  \item Nepregledno - 5 razina prosljeđivanja za jedan podatak
  \item Krto - promjena na jednoj razini razbija lanac
  \item Teško za maintain - ne znaš odakle podatak dolazi
\end{itemize}

\vspace{0.6em}
Store rješava ovo: bilo koja komponenta može direktno čitati i pisati dijeljeni state bez posrednika.
\end{frame}

% ------------------------------------------------------------
\begin{frame}{Pinia vs Vuex}
\begin{tabular}{@{}lp{5cm}p{5cm}@{}}
  \toprule
  & \textbf{Vuex (legacy)} & \textbf{Pinia} \\
  \midrule
  API & Options: state, mutations, actions, getters & Composition: ref, computed, function \\
  Mutacije & Obavezne (sinkrone) & Nema - direktno mijenjaj state \\
  TypeScript & Loša podrška & Odlična \\
  DevTools & Da & Da (bolje) \\
  Moduli & Kompleksni namespace & Svaki store je nezavisan \\
  \bottomrule
\end{tabular}

\vspace{0.6em}
Pinia je službeni nasljednik Vuexa. Ako naiđete na tutoriale s Vuexom, koncepti su isti, ali API je drugačiji.
\end{frame}

% ------------------------------------------------------------
\begin{frame}{Store pattern: state, actions, getters}
Svaki Pinia store ima tri koncepta:

\vspace{0.4em}
\begin{tabular}{@{}lp{4cm}p{5cm}@{}}
  \toprule
  \textbf{Koncept} & \textbf{U Piniji} & \textbf{Analogija (backend)} \\
  \midrule
  State & \texttt{ref()} & Podatak u bazi \\
  Getters & \texttt{computed()} & Izvedeni podatak (SQL VIEW) \\
  Actions & Funkcije & Service metode \\
  \bottomrule
\end{tabular}

\vspace{0.6em}
\begin{itemize}
  \item \textbf{State} drži sirove podatke (lista natjecatelja, loading flag)
  \item \textbf{Getters} deriviraju podatke iz statea (filtrirani natjecatelji, ukupan broj)
  \item \textbf{Actions} mijenjaju state (dohvat s API-ja, dodavanje, brisanje)
\end{itemize}
\end{frame}

% ------------------------------------------------------------
\begin{frame}{Optimistički vs pesimistički update}
\begin{columns}[T]
\column{0.48\textwidth}
\textbf{Pesimistički}
\begin{itemize}
  \item Pošalji request $\rightarrow$ čekaj odgovor $\rightarrow$ ažuriraj UI
  \item Korisnik vidi spinner
  \item Sigurniji - nema rollbacka
  \item Sporiji UX
\end{itemize}

\column{0.48\textwidth}
\textbf{Optimistički}
\begin{itemize}
  \item Ažuriraj UI odmah $\rightarrow$ pošalji request $\rightarrow$ rollback ako ne uspije
  \item Korisnik vidi instant promjenu
  \item Kompleksniji - treba rollback logika
  \item Brži UX
\end{itemize}
\end{columns}

\vspace{0.6em}
Za ovaj projekt koristimo pesimistički pristup - jednostavniji je i dovoljno brz. Optimistički update se koristi za akcije poput lajkanja objave ili micanja itema iz liste.
\end{frame}

% ------------------------------------------------------------
\begin{frame}[fragile,shrink=5]{Pinia store za natjecatelje}
\begin{lstlisting}[language=typescript]
// src/stores/lifters.ts
import { defineStore } from 'pinia'
import { ref } from 'vue'
import api from '@/api/client'
import type { Lifter } from '@/types'

export const useLifterStore = defineStore('lifters', () => {
  const lifters = ref<Lifter[]>([])
  const loading = ref(false)
  const error = ref<string | null>(null)

  async function fetchLifters(clubId: number) {
    loading.value = true
    error.value = null
    try {
      const { data } = await api.get<Lifter[]>(
        `/clubs/${clubId}/lifters`
      )
      lifters.value = data
    } catch {
      error.value = 'Greska pri dohvatu natjecatelja'
    } finally {
      loading.value = false
    }
  }

  return { lifters, loading, error, fetchLifters }
})
\end{lstlisting}
\end{frame}

% ------------------------------------------------------------
\begin{frame}[fragile,shrink=5]{Zamka: destrukturiranje storea gubi reaktivnost}
\begin{lstlisting}[language=typescript]
const store = useLifterStore()
const { lifters } = store          // GUBI reaktivnost!

import { storeToRefs } from 'pinia'
const { lifters, loading } = storeToRefs(store)  // ispravno
const { fetchLifters } = store     // funkcije smiju direktno
\end{lstlisting}

\vspace{0.4em}
Destrukturiranje kopira trenutnu vrijednost i template se prestaje ažurirati kad se store promijeni. \texttt{storeToRefs} zadržava reaktivne reference, a akcije se smiju uzeti direktno jer su obične funkcije.

\vspace{0.4em}
\textbf{Napomena:} AI alati često generiraju upravo ovu grešku - provjerite svako destrukturiranje storea koje vam alat predloži.
\end{frame}

% ------------------------------------------------------------
\begin{frame}{Tri stanja svakog viewa}
Svaki view koji dohvaća podatke s API-ja mora hendlati tri situacije:

\vspace{0.6em}
\begin{tabular}{@{}llp{6cm}@{}}
  \toprule
  \textbf{Stanje} & \textbf{Uvjet} & \textbf{Što prikazujemo} \\
  \midrule
  Loading & \texttt{loading === true} & Spinner ili skeleton \\
  Error & \texttt{error !== null} & Poruka greške s retry gumbom \\
  Empty & \texttt{items.length === 0} & Poruka i CTA (npr. \enquote{Dodaj prvog natjecatelja}) \\
  \bottomrule
\end{tabular}

\vspace{0.6em}
Ako pokažete samo podatke bez ovih stanja, korisnik ne zna učitava li se nešto, je li greška ili stvarno nema podataka. To je loš UX.
\end{frame}

% ------------------------------------------------------------
\begin{frame}{Skeleton loading vs spinner}
\begin{columns}[T]
\column{0.48\textwidth}
\textbf{Spinner}
\begin{itemize}
  \item Generički indikator učitavanja
  \item Ne govori korisniku što će vidjeti
  \item Jednostavan za implementaciju
  \item Može uzrokovati layout shift
\end{itemize}

\column{0.48\textwidth}
\textbf{Skeleton screen}
\begin{itemize}
  \item Sivi blokovi na mjestu budućeg sadržaja
  \item Korisnik zna kakav layout očekivati
  \item Perceptualno brže (manje čekanja)
  \item Zahtijeva više koda
\end{itemize}
\end{columns}

\vspace{0.6em}
Za studentski projekt spinner je dovoljno dobar. Skeleton je standard u produkcijskim aplikacijama - Facebook, YouTube, LinkedIn ga koriste svugdje.
\end{frame}

% ------------------------------------------------------------
\begin{frame}{Empty state: zašto je bitan}
Empty state je ono što korisnik vidi kad nema podataka. Loš empty state:

\vspace{0.2em}
\begin{center}
\textit{(prazna tablica bez ikakvog teksta)}
\end{center}

\vspace{0.4em}
Dobar empty state:
\begin{itemize}
  \item Jasna poruka: \enquote{Nemate natjecatelja.}
  \item CTA (Call to Action): link ili gumb za dodavanje prvog zapisa
  \item Opcionalno: kratka uputa što sljedeće napraviti
\end{itemize}

\vspace{0.6em}
Empty state se događa češće nego mislite - svaki novi korisnik vidi ga prvi put. To je prilika za dobar prvi dojam, ne prazna stranica.
\end{frame}

% ------------------------------------------------------------
\begin{frame}{Error stanja i retry}
Nije svaka greška ista. Razlikujemo:

\vspace{0.4em}
\begin{tabular}{@{}lp{4cm}p{5cm}@{}}
  \toprule
  \textbf{Tip} & \textbf{Primjer} & \textbf{Reakcija} \\
  \midrule
  Network error & Nema interneta & Retry gumb, provjeri konekciju \\
  Server error (5xx) & Bug na backendu & Retry gumb, \enquote{Pokušajte kasnije} \\
  Auth error (401) & Token istekao & Automatski refresh ili redirect \\
  Client error (4xx) & Krivi podaci & Prikaži poruku, ne nudi retry \\
  \bottomrule
\end{tabular}

\vspace{0.6em}
Retry ima smisla samo za tranzijentne greške (network, 5xx). Za 400/422/409 retry će dati isti rezultat - korisnik mora promijeniti input.
\end{frame}

% ------------------------------------------------------------
\begin{frame}{Debouncing: ne šalji request na svaki keystroke}
Kad korisnik piše u search polje, ne želimo slati API request na svako slovo:

\vspace{0.4em}
\begin{center}
\texttt{p $\rightarrow$ po $\rightarrow$ pow $\rightarrow$ powe $\rightarrow$ power} = 5 requestova za jednu riječ
\end{center}

\vspace{0.6em}
\textbf{Debounce} odgađa akciju dok korisnik ne prestane tipkati (npr. 300ms):
\begin{enumerate}
  \item Korisnik piše \texttt{p} $\rightarrow$ timer kreće (300ms)
  \item Piše \texttt{o} $\rightarrow$ timer se resetira
  \item Piše \texttt{w}, \texttt{e}, \texttt{r} $\rightarrow$ timer se resetira svaki put
  \item 300ms bez tipkanja $\rightarrow$ šalje se jedan request za \texttt{power}
\end{enumerate}

\vspace{0.4em}
Za frontend filtriranje (computed) debounce nije potreban jer nema network poziva. Za API search endpoint je obavezan.
\end{frame}

% ------------------------------------------------------------
\begin{frame}[fragile,shrink=5]{Composable: useDebouncedSearch}
\begin{lstlisting}[language=typescript]
// src/composables/useDebouncedSearch.ts
import { ref, watch } from 'vue'

export function useDebouncedSearch(delay = 300) {
  const search = ref('')
  const debounced = ref('')
  let timer: ReturnType<typeof setTimeout>

  watch(search, (value) => {
    clearTimeout(timer)
    timer = setTimeout(() => { debounced.value = value }, delay)
  })

  return { search, debounced }
}

// u komponenti:
// const { search, debounced } = useDebouncedSearch()
// <input v-model="search" />  -> API poziv prati debounced
\end{lstlisting}

\vspace{0.4em}
Composable je funkcija s prefiksom \texttt{use} koja kapsulira reaktivnu logiku - isto što i Pinia store, ali bez globalnog stanja. Svaka komponenta koja je pozove dobiva svoju instancu. Ovo je centralni Composition API pattern za dijeljenje logike.
\end{frame}

% ------------------------------------------------------------
\begin{frame}[fragile,shrink=5]{Lista natjecatelja s UX stanjima}
\begin{lstlisting}[language=html]
<script setup lang="ts">
import { onMounted } from 'vue'
import { useAuthStore } from '@/stores/auth'
import { useLifterStore } from '@/stores/lifters'

const auth = useAuthStore()
const store = useLifterStore()

onMounted(() => {
  store.fetchLifters(auth.user.club_id)
})
</script>

<template>
  <h1>Natjecatelji</h1>
  <div v-if="store.loading">Ucitavanje...</div>
  <div v-else-if="store.error" class="error">
    {{ store.error }}
    <button @click="store.fetchLifters(auth.user.club_id)">
      Pokusaj ponovo
    </button>
  </div>
  <div v-else-if="store.lifters.length === 0">
    Nemate natjecatelja. <RouterLink to="/lifters/new">
    Dodajte prvog.</RouterLink>
  </div>
  <table v-else>...</table>
</template>
\end{lstlisting}
\end{frame}

% ------------------------------------------------------------
\begin{frame}[fragile,shrink=5]{Tablica podataka}
\begin{lstlisting}[language=html]
<table v-else>
  <thead>
    <tr>
      <th>Ime</th>
      <th>Prezime</th>
      <th>Spol</th>
      <th>Datum rodenja</th>
      <th>Akcije</th>
    </tr>
  </thead>
  <tbody>
    <tr v-for="lifter in store.lifters" :key="lifter.id">
      <td>{{ lifter.first_name }}</td>
      <td>{{ lifter.last_name }}</td>
      <td>{{ lifter.gender }}</td>
      <td>{{ lifter.birth_date }}</td>
      <td>
        <RouterLink :to="`/lifters/${lifter.id}/edit`">
          Uredi
        </RouterLink>
      </td>
    </tr>
  </tbody>
</table>
\end{lstlisting}

\vspace{0.4em}
\texttt{v-for} iterira kroz array. \texttt{:key} je obavezan - Vue ga koristi za efikasno ažuriranje DOM-a.
\end{frame}

% ------------------------------------------------------------
\begin{frame}[fragile,shrink=5]{Forma za dodavanje natjecatelja}
\begin{lstlisting}[language=html]
<script setup lang="ts">
import { ref } from 'vue'
import { useRouter } from 'vue-router'
import { AxiosError } from 'axios'
import api from '@/api/client'
import { useAuthStore } from '@/stores/auth'

const auth = useAuthStore()
const router = useRouter()
const form = ref({
  first_name: '', last_name: '',
  gender: 'M', birth_date: '',
})
const errors = ref<Record<string, string>>({})

async function submit() {
  errors.value = {}
  try {
    await api.post(
      `/clubs/${auth.user!.club_id}/lifters`, form.value
    )
    router.push('/lifters')
  } catch (e) {
    const err = e as AxiosError
    if (err.response?.status === 422) {
      errors.value = parseValidationErrors(err.response.data)
    }
  }
}
</script>
\end{lstlisting}
\end{frame}

% ------------------------------------------------------------
\begin{frame}[fragile,shrink=5]{Forma: template s validacijom}
\begin{lstlisting}[language=html]
<template>
  <h1>Novi natjecatelj</h1>
  <form @submit.prevent="submit">
    <div>
      <label>Ime</label>
      <input v-model="form.first_name" required />
      <span v-if="errors.first_name" class="error">
        {{ errors.first_name }}
      </span>
    </div>
    <div>
      <label>Spol</label>
      <select v-model="form.gender">
        <option value="M">Muski</option>
        <option value="F">Zenski</option>
      </select>
    </div>
    <div>
      <label>Datum rodenja</label>
      <input v-model="form.birth_date" type="date" required />
    </div>
    <button type="submit">Spremi</button>
  </form>
</template>
\end{lstlisting}
\end{frame}

% ------------------------------------------------------------
\begin{frame}[fragile,shrink=5]{Edit forma: punjenje postojećim podacima}
\begin{lstlisting}[language=html]
// src/views/LifterEditView.vue
<script setup lang="ts">
import { ref, onMounted, watch } from 'vue'
import { useRoute } from 'vue-router'
import api from '@/api/client'
import type { Lifter } from '@/types'

const route = useRoute()
const form = ref({ first_name: '', last_name: '',
                   gender: 'M', birth_date: '' })

async function load(id: number) {
  const { data } = await api.get<Lifter>(`/lifters/${id}`)
  form.value = { ...data }
}

onMounted(() => load(Number(route.params.id)))
watch(() => route.params.id,
      (id) => load(Number(id)))
</script>
\end{lstlisting}

\vspace{0.4em}
\texttt{onMounted} pokriva prvi ulazak na rutu, a \texttt{watch} na route param slučaj kad router ponovno iskoristi istu komponentu za drugi ID - bez watcha forma ostane s podacima starog zapisa. Ovo je najčešći praktični \texttt{watch} u SPA aplikaciji.
\end{frame}

% ------------------------------------------------------------
\begin{frame}{Validacija: klijentska vs serverska}
\begin{columns}[T]
\column{0.48\textwidth}
\textbf{Klijentska (frontend)}
\begin{itemize}
  \item HTML atributi: \texttt{required}, \texttt{min}, \texttt{maxlength}
  \item Brzi feedback - bez mrežnog poziva
  \item Može se zaobići (DevTools)
  \item Za UX, ne za sigurnost
\end{itemize}

\column{0.48\textwidth}
\textbf{Serverska (backend)}
\begin{itemize}
  \item Pydantic schema validacija
  \item Vraća 422 s detaljima
  \item Ne može se zaobići
  \item Autoritet za podatke
\end{itemize}
\end{columns}

\vspace{0.6em}
Trebate obje. Klijentska daje brzi feedback, serverska jamči integritet. Frontend prikazuje serverske greške kad klijentska validacija nije dovoljna (npr. duplikat u bazi).
\end{frame}

% ------------------------------------------------------------
\begin{frame}{Form UX best practices}
\begin{enumerate}
  \item \textbf{Inline validacija} - prikaži grešku ispod polja odmah kad korisnik napusti polje, ne tek na submit
  \item \textbf{Disable submit dok forma nije validna} - korisnik zna da nešto nedostaje
  \item \textbf{Clear error kad korisnik počne tipkati} - ne drži crveni tekst nakon što je korisnik počeo ispravljati
  \item \textbf{Focus na prvo polje s greškom} - nakon neuspjelog submita, premjesti fokus na problem
  \item \textbf{Serverske greške uz polje} - 422 greška za \texttt{birth\_date} ide pored datumskog inputa, ne kao generička poruka
\end{enumerate}
\end{frame}

% ------------------------------------------------------------
\begin{frame}{Prikaz serverskih grešaka u formi}
Backend može vratiti dva tipa grešaka koje trebamo prikazati:

\vspace{0.4em}
\begin{tabular}{@{}llp{5.5cm}@{}}
  \toprule
  \textbf{Status} & \textbf{Tip} & \textbf{Prikaz} \\
  \midrule
  422 & Field-level validacija & Greška ispod odgovarajućeg polja (\texttt{loc} kaže koje) \\
  409 & Duplikat & Globalna poruka iznad forme: \enquote{Natjecatelj je već prijavljen} \\
  400 & Poslovno pravilo & Globalna poruka: \enquote{Rok za prijave je istekao} \\
  \bottomrule
\end{tabular}

\vspace{0.6em}
Razlikujemo ih po status kodu. 422 ima \texttt{detail} array s lokacijama polja, 409 i 400 imaju \texttt{message} string. Frontend parsira response i prikazuje grešku na pravom mjestu.
\end{frame}

% ------------------------------------------------------------
\begin{frame}[fragile,shrink=5]{Parsiranje 422 grešaka s backenda}
\begin{lstlisting}[language=typescript]
// src/utils/errors.ts

// FastAPI 422 format:
// { "detail": [{ "loc": ["body","first_name"],
//   "msg": "field required", "type": "missing" }] }

interface ValidationDetail {
  loc: (string | number)[]
  msg: string
}

export function parseValidationErrors(
  data: { detail?: ValidationDetail[] } | unknown
): Record<string, string> {
  const errors: Record<string, string> = {}
  const detail = (data as { detail?: ValidationDetail[] })?.detail
  if (Array.isArray(detail)) {
    for (const err of detail) {
      // loc je ["body", "field_name"]
      const field = String(err.loc[err.loc.length - 1])
      errors[field] = err.msg
    }
  }
  return errors
}
\end{lstlisting}

\vspace{0.4em}
Svaka greška ima \texttt{loc} (lokacija polja) i \texttt{msg} (poruka). Mapiramo ih na objekt \texttt{\{field: message\}} za prikaz u formi.
\end{frame}

% ------------------------------------------------------------
\begin{frame}[fragile,shrink=5]{Registracije: prikaz sa status badgeom}
\begin{lstlisting}[language=html]
<template>
  <table>
    <tr v-for="reg in registrations" :key="reg.id">
      <td>{{ reg.lifter_name }}</td>
      <td>{{ reg.category }}</td>
      <td>
        <span :class="statusClass(reg.status)">
          {{ reg.status }}
        </span>
      </td>
      <td>
        <button v-if="reg.status === 'active'"
                @click="withdraw(reg.id)">
          Odjavi
        </button>
      </td>
    </tr>
  </table>
</template>

<script setup lang="ts">
function statusClass(status: string): string {
  return status === 'active' ? 'badge-green' : 'badge-red'
}
</script>
\end{lstlisting}
\end{frame}

% ------------------------------------------------------------
\begin{frame}{Reusable komponente: kad izdvojiti}
Izdvajamo komponentu kad se isti UI pattern ponavlja na 2+ mjesta:

\vspace{0.4em}
\begin{tabular}{@{}ll@{}}
  \toprule
  \textbf{Komponenta} & \textbf{Razlog} \\
  \midrule
  \codefile{LoadingSpinner.vue} & Koristi se u svakom viewu \\
  \codefile{ErrorMessage.vue} & Isti pattern: poruka + retry \\
  \codefile{DataTable.vue} & Tablica s headerom i redovima \\
  \codefile{FormField.vue} & Label + input + error poruka \\
  \codefile{StatusBadge.vue} & Obojeni badge za statuse \\
  \bottomrule
\end{tabular}

\vspace{0.6em}
Ne pretjerujte s apstrakcijom. Ako se nešto koristi samo jednom, ne treba biti zasebna komponenta.
\end{frame}

% ------------------------------------------------------------
\begin{frame}{Kompozicija vs nasljeđivanje u UI-ju}
Vue ne koristi nasljeđivanje komponenti. Umjesto toga koristi \textbf{kompoziciju}:

\vspace{0.4em}
\begin{tabular}{@{}lp{5cm}p{5cm}@{}}
  \toprule
  & \textbf{Nasljeđivanje} & \textbf{Kompozicija} \\
  \midrule
  Pristup & \texttt{class AdminTable extends Table} & \texttt{<DataTable>} s propsima i slotovima \\
  Fleksibilnost & Rigidna hijerarhija & Slobodna kombinacija \\
  Reusability & Teško za različite kontekste & Lako - props i slots \\
  \bottomrule
\end{tabular}

\vspace{0.6em}
Slots omogućuju parentu da definira \textbf{što} se prikazuje, a child komponenta definira \textbf{gdje} i \textbf{kako}. Ovo je isti princip kao \enquote{composition over inheritance} iz OOP-a.
\end{frame}

% ------------------------------------------------------------
\begin{frame}{Props dizajn: API vaše komponente}
Propsi su \textbf{javni API} vaše komponente. Dizajnirajte ih s istom pažnjom kao REST endpoint:

\vspace{0.4em}
\begin{itemize}
  \item \textbf{Imenujte opisno}: \texttt{isLoading} umjesto \texttt{flag}, \texttt{errorMessage} umjesto \texttt{msg}
  \item \textbf{Definirajte tipove}: generički \texttt{defineProps<\{...\}>()} - TypeScript provjerava svaki prop na mjestu korištenja
  \item \textbf{Postavite defaulte}: \texttt{withDefaults()} za opcionalne propse
  \item \textbf{Koristite v-model za dvosmjerne propse}: umjesto prop + emit pattern ručno
\end{itemize}

\vspace{0.6em}
Komponenta s 10+ propsova je znak da radi previše. Razmislite o razbijanju na manje komponente ili korištenju slotova.
\end{frame}

% ------------------------------------------------------------
\begin{frame}{Key prop u v-for: zašto je obavezan}
Bez \texttt{:key}, Vue koristi \enquote{in-place patch} strategiju - ažurira postojeće DOM elemente umjesto da ih premjesti.

\vspace{0.4em}
Problem: korisnik uređuje input u 2. redu tablice. Filtriranjem se 2. red promijeni - ali DOM element ostane isti. Input vrijednost pripada starom redu.

\vspace{0.6em}
\begin{tabular}{@{}lp{8cm}@{}}
  \toprule
  \textbf{Key} & \textbf{Ponašanje} \\
  \midrule
  \texttt{:key="item.id"} & Vue prati identitet elementa - premješta DOM node kad se lista presloži \\
  \texttt{:key="index"} & Vue tretira poziciju kao identitet - bugovi pri sortiranju/filtriranju \\
  Bez keya & Warning u konzoli, nepredvidljivo ponašanje \\
  \bottomrule
\end{tabular}

\vspace{0.4em}
Koristite jedinstven, stabilan identifikator (ID iz baze) kao key. Nikad index.
\end{frame}

% ------------------------------------------------------------
\begin{frame}{Emit pattern: inverzija kontrole}
Child komponenta ne odlučuje što se događa kad korisnik klikne gumb - samo javlja parentu.

\vspace{0.6em}
\begin{center}
\texttt{Child: emit('delete', itemId)} $\rightarrow$ \texttt{Parent: handleDelete(itemId)}
\end{center}

\vspace{0.6em}
\begin{itemize}
  \item Child ne zna koristi li se u admin viewu ili club viewu
  \item Parent odlučuje: prikaži modal? Odmah obriši? Zatraži razlog?
  \item Ista child komponenta se koristi u različitim kontekstima s različitim ponašanjem
\end{itemize}

\vspace{0.6em}
Ovo je \textbf{inverzija kontrole} - isti princip kao dependency injection na backendu. Komponenta prima ponašanje izvana umjesto da ga hardkodira.
\end{frame}

% ------------------------------------------------------------
\begin{frame}[fragile,shrink=5]{Reusable FormField komponenta}
\begin{lstlisting}[language=html]
<!-- src/components/FormField.vue -->
<script setup lang="ts">
defineProps<{
  label: string
  error?: string
  modelValue: string | number
}>()
defineEmits<{
  'update:modelValue': [value: string]
}>()
</script>

<template>
  <div class="form-field">
    <label>{{ label }}</label>
    <input :value="modelValue"
      @input="$emit('update:modelValue',
        ($event.target as HTMLInputElement).value)"
    />
    <span v-if="error" class="error">{{ error }}</span>
  </div>
</template>

<!-- Koristenje: -->
<!-- <FormField v-model="form.first_name" label="Ime"
     :error="errors.first_name" /> -->
\end{lstlisting}
\end{frame}

% ------------------------------------------------------------
\begin{frame}{Commitovi za Predavanje 9}
\begin{itemize}
  \item \texttt{feat(web): add lifters list view with loading/error/empty states}
  \item \texttt{feat(web): add lifter create/edit form with validation}
  \item \texttt{feat(web): add competitions list and registration form}
  \item \texttt{feat(web): add my registrations view with status badges}
  \item \texttt{refactor(web): extract reusable form and table components}
\end{itemize}

\vspace{0.4em}
\textbf{Na satu:} gradimo CRUD viewove, testiramo forme, promatramo Network tab u DevToolsima.
\end{frame}

% ------------------------------------------------------------
\begin{frame}{Što ide sljedeće (Predavanje 10)}
\begin{itemize}
  \item Admin viewovi: tablice s filterima i paginacijom
  \item Permission-aware UI - prikaz/sakrivanje ovisno o roli i fazi
  \item Modali za potvrdu akcija
  \item CSV export
\end{itemize}

\vspace{0.8em}
\begin{center}
\Large Pitanja?
\end{center}
\end{frame}

\end{document}
